% ---- score2 standard preamble (from: python3 score2_tools.py preamble) ----
\documentclass[10pt]{extarticle}
\usepackage[a4paper,margin=1.5cm]{geometry}
\usepackage[T1]{fontenc}
\usepackage{amsmath,amssymb}
\usepackage{graphicx}
\usepackage[export]{adjustbox}
\usepackage{multicol}
\usepackage{enumitem}
\usepackage{booktabs,array,tabularx}
\setlength{\columnsep}{0.7cm}
\setlength{\parindent}{0pt}
\setlength{\parskip}{3pt}
\pagestyle{empty}
\makeatletter
% figure registry (figdefs.tex is generated by the extract step)
\newcommand{\qfigdef}[4]{\@namedef{qf@file@#1}{#2}\@namedef{qf@w@#1}{#3}\@namedef{qf@wide@#1}{#4}}
\newcommand{\qf@check}[1]{\@ifundefined{qf@file@#1}{\PackageError{score2}{Unknown figure id #1}{}}{}}
\newcommand{\qf@img}[4][]{\includegraphics[width=#2\dimexpr\@nameuse{qf@w@#4}pt\relax,max width=#3,max height=0.45\textheight,#1]{\@nameuse{qf@file@#4}}}
\newif\ifqf@full
\newif\ifqf@boxed
\newsavebox{\qf@box}
\newcommand{\qf@open}{\global\qf@fulltrue\par\global\qf@boxedtrue
  \setbox\qf@box\vbox\bgroup\hsize\textwidth\linewidth\textwidth}
\newcommand{\qf@emit}[1]{\par\egroup\global\qf@boxedfalse
  \ifdim\dimexpr\ht\qf@box+\dp\qf@box\relax<#1\textheight\noindent\usebox{\qf@box}\par\else\unvbox\qf@box\fi}
\newcommand{\qfullsplit}{\ifqf@boxed\qf@emit{0.85}\fi}
\newcommand{\qf@close}{\ifqf@boxed\qf@emit{0.45}\fi\global\qf@fullfalse}
\newcommand{\qfullbegin}{\end{multicols}\qf@open}
\newcommand{\qfullend}{\qf@close\begin{multicols}{2}}
\newcommand{\qfullstart}{\qf@open}
\newcommand{\qfulldone}{\qf@close}
\newcommand{\qfig}[2][1]{\qf@check{#2}%
  \ifqf@full\qf@blk{#1}{#2}\else
  \ifnum\@nameuse{qf@wide@#2}=1 \expandafter\@firstoftwo\else\expandafter\@secondoftwo\fi
  {\end{multicols}\begin{center}\qf@img{#1}{\textwidth}{#2}\end{center}\begin{multicols}{2}}%
  {\qf@blk{#1}{#2}}\fi}
\newcommand{\qf@blk}[2]{\par\smallskip{\centering\qf@img{#1}{\linewidth}{#2}\par}\smallskip}
\newcommand{\qfigi}[2][1]{\qf@check{#2}\qf@img[valign=c]{#1}{\linewidth}{#2}}
\makeatother
\newcommand{\papersec}[1]{\par\bigskip{\centering\Large\bfseries PAPER-#1\par}\medskip}
\newcommand{\partsec}[2]{\par\medskip\noindent{\large\bfseries PART-#1: #2}\par\smallskip}
\newcommand{\sect}[1]{\par\smallskip\noindent\textbf{#1}\par}
\newcommand{\qn}[1]{\par\smallskip\noindent\textbf{#1.}~}
\newcommand{\paraq}[2]{\par\smallskip\noindent\textbf{Paragraph for Questions #1 and #2}\par}
\newenvironment{qoptions}{\begin{itemize}[leftmargin=1.8em,itemsep=1pt,topsep=2pt,label={}]}{\end{itemize}}
\newcommand{\opt}[2]{\item[(#1)] #2}
% figure registry (generated by score2_tools.py finalize)
\qfigdef{p1_phy_s1_q01_1}{Score-II - Test 09 (12 Feb 2023)-img/p1_phy_s1_q01_1.png}{108.2}{0}
\qfigdef{p1_phy_s1_q02_1}{Score-II - Test 09 (12 Feb 2023)-img/p1_phy_s1_q02_1.png}{100.4}{0}
\qfigdef{p1_phy_s1_q03_1}{Score-II - Test 09 (12 Feb 2023)-img/p1_phy_s1_q03_1.png}{108.0}{0}
\qfigdef{p1_phy_s1_q04_1}{Score-II - Test 09 (12 Feb 2023)-img/p1_phy_s1_q04_1.png}{94.8}{0}
\qfigdef{p1_phy_s1_q05_1}{Score-II - Test 09 (12 Feb 2023)-img/p1_phy_s1_q05_1.png}{97.2}{0}
\qfigdef{p1_phy_s1_q06_1}{Score-II - Test 09 (12 Feb 2023)-img/p1_phy_s1_q06_1.png}{95.5}{0}
\qfigdef{p1_phy_s2_q07_1}{Score-II - Test 09 (12 Feb 2023)-img/p1_phy_s2_q07_1.png}{89.2}{0}
\qfigdef{p1_phy_s2_q08_1}{Score-II - Test 09 (12 Feb 2023)-img/p1_phy_s2_q08_1.png}{140.1}{0}
\qfigdef{p1_phy_s2_q09_1}{Score-II - Test 09 (12 Feb 2023)-img/p1_phy_s2_q09_1.png}{132.5}{0}
\qfigdef{p1_phy_s2_q10_1}{Score-II - Test 09 (12 Feb 2023)-img/p1_phy_s2_q10_1.png}{125.8}{0}
\qfigdef{p1_phy_s3_q03_1}{Score-II - Test 09 (12 Feb 2023)-img/p1_phy_s3_q03_1.png}{160.6}{0}
\qfigdef{p1_phy_s3_q04_1}{Score-II - Test 09 (12 Feb 2023)-img/p1_phy_s3_q04_1.png}{174.1}{0}
\qfigdef{p1_phy_s3_q05_1}{Score-II - Test 09 (12 Feb 2023)-img/p1_phy_s3_q05_1.png}{132.8}{0}
\qfigdef{p1_chem_s1_q01_1}{Score-II - Test 09 (12 Feb 2023)-img/p1_chem_s1_q01_1.png}{208.9}{0}
\qfigdef{p1_chem_s1_q03_1}{Score-II - Test 09 (12 Feb 2023)-img/p1_chem_s1_q03_1.png}{164.9}{0}
\qfigdef{p1_chem_s1_q03_2}{Score-II - Test 09 (12 Feb 2023)-img/p1_chem_s1_q03_2.png}{129.6}{0}
\qfigdef{p1_chem_s1_q03_3}{Score-II - Test 09 (12 Feb 2023)-img/p1_chem_s1_q03_3.png}{76.8}{0}
\qfigdef{p1_chem_s1_q03_4}{Score-II - Test 09 (12 Feb 2023)-img/p1_chem_s1_q03_4.png}{176.1}{0}
\qfigdef{p1_chem_s1_q05_1}{Score-II - Test 09 (12 Feb 2023)-img/p1_chem_s1_q05_1.png}{85.2}{0}
\qfigdef{p1_chem_s2_q08_1}{Score-II - Test 09 (12 Feb 2023)-img/p1_chem_s2_q08_1.png}{97.9}{0}
\qfigdef{p1_chem_s3_q02_1}{Score-II - Test 09 (12 Feb 2023)-img/p1_chem_s3_q02_1.png}{71.1}{0}
\qfigdef{p1_chem_s3_q05_1}{Score-II - Test 09 (12 Feb 2023)-img/p1_chem_s3_q05_1.png}{243.2}{0}
\qfigdef{p2_phy_s1_q01_1}{Score-II - Test 09 (12 Feb 2023)-img/p2_phy_s1_q01_1.png}{90.1}{0}
\qfigdef{p2_phy_s1_q02_1}{Score-II - Test 09 (12 Feb 2023)-img/p2_phy_s1_q02_1.png}{112.3}{0}
\qfigdef{p2_phy_s1_q04_1}{Score-II - Test 09 (12 Feb 2023)-img/p2_phy_s1_q04_1.png}{68.9}{0}
\qfigdef{p2_phy_s2_q05_1}{Score-II - Test 09 (12 Feb 2023)-img/p2_phy_s2_q05_1.png}{213.9}{0}
\qfigdef{p2_phy_s2_q06_1}{Score-II - Test 09 (12 Feb 2023)-img/p2_phy_s2_q06_1.png}{104.9}{0}
\qfigdef{p2_phy_s2_q07_1}{Score-II - Test 09 (12 Feb 2023)-img/p2_phy_s2_q07_1.png}{123.4}{0}
\qfigdef{p2_phy_s2_q08_1}{Score-II - Test 09 (12 Feb 2023)-img/p2_phy_s2_q08_1.png}{67.3}{0}
\qfigdef{p2_phy_s3_q02_1}{Score-II - Test 09 (12 Feb 2023)-img/p2_phy_s3_q02_1.png}{89.7}{0}
\qfigdef{p2_phy_s3_q03_1}{Score-II - Test 09 (12 Feb 2023)-img/p2_phy_s3_q03_1.png}{87.7}{0}
\qfigdef{p2_phy_s3_q05_1}{Score-II - Test 09 (12 Feb 2023)-img/p2_phy_s3_q05_1.png}{122.4}{0}
\qfigdef{p2_phy_s3_q06_1}{Score-II - Test 09 (12 Feb 2023)-img/p2_phy_s3_q06_1.png}{78.7}{0}
\qfigdef{p2_phy_s3_q07_1}{Score-II - Test 09 (12 Feb 2023)-img/p2_phy_s3_q07_1.png}{63.5}{0}
\qfigdef{p2_phy_s3_q08_1}{Score-II - Test 09 (12 Feb 2023)-img/p2_phy_s3_q08_1.png}{118.3}{0}
\qfigdef{p2_chem_s1_q01_1}{Score-II - Test 09 (12 Feb 2023)-img/p2_chem_s1_q01_1.png}{159.2}{0}
\qfigdef{p2_chem_s2_q08_1}{Score-II - Test 09 (12 Feb 2023)-img/p2_chem_s2_q08_1.png}{267.7}{1}
\qfigdef{p2_chem_s2_q08_2}{Score-II - Test 09 (12 Feb 2023)-img/p2_chem_s2_q08_2.png}{152.5}{0}
\qfigdef{p2_chem_s2_q08_3}{Score-II - Test 09 (12 Feb 2023)-img/p2_chem_s2_q08_3.png}{72.3}{0}
\qfigdef{p2_chem_s2_q09_1}{Score-II - Test 09 (12 Feb 2023)-img/p2_chem_s2_q09_1.png}{160.4}{0}
\qfigdef{p2_chem_s2_q09_2}{Score-II - Test 09 (12 Feb 2023)-img/p2_chem_s2_q09_2.png}{47.7}{0}
\qfigdef{p2_chem_s2_q09_3}{Score-II - Test 09 (12 Feb 2023)-img/p2_chem_s2_q09_3.png}{55.8}{0}
\qfigdef{p2_chem_s3_q03_1}{Score-II - Test 09 (12 Feb 2023)-img/p2_chem_s3_q03_1.png}{155.4}{0}
\begin{document}
\begin{center}{\Large\bfseries JEE (Main+Advanced) Score-II --- Test 09 (12 Feb 2023)}\end{center}
\begin{multicols}{2}
\papersec{1}
\partsec{1}{PHYSICS}
\sect{Section-I (i): One or More Correct}

\qn{1} A unimolar sample of argon is undergoing a cyclic process as indicated by a T--V indicator diagram shown. Which of the following alternatives is/are correct?
\qfig{p1_phy_s1_q01_1}
\begin{qoptions}
\opt{A}{During the process DA the pressure of the gas increases from $P_D$ to $P_A$ and is equal to $P_C$ sometime during the process.}
\opt{B}{$|Q_{AB}|>|Q_{BC}|=|Q_{DA}|>|Q_{CD}|$}
\opt{C}{Efficiency of the cycle is $\dfrac{2}{5+4\ln 2}$.}
\opt{D}{Maximum pressure during the cycle is $\dfrac{RT_0}{V_0}$.}
\end{qoptions}

\qn{2} In the arrangement shown there are four identical capacitors of capacitance C connected to a cell of emf $\varepsilon$. Two attempts are made to change the capacitance of the circuit. In the first attempt the capacitor B is filled completely with dielectric of dielectric constant 3 and in the second attempt in the original circuit capacitor A is filled completely with the same dielectric. Choose the correct alternatives, considering that circuit aquires steady state after the changes.
\qfig{p1_phy_s1_q02_1}
\begin{qoptions}
\opt{A}{The ratio of change in equivalent capacitance in attempt 1 to attempt 2 is $\dfrac{11}{36}$.}
\opt{B}{The ratio of the work done by the cell in attempt 1 and attempt 2 is $\dfrac{11}{36}$.}
\opt{C}{The ratio of energy stored in capacitor A and B after attempt-1 will be $16:3$.}
\opt{D}{The amount of heat dissipated in attempt-1 is $\dfrac{1}{45}C\varepsilon^2$.}
\end{qoptions}

\qn{3} Two rings each of mass m and radius R are rigidly attached to each other such that their planes are perpendicular to each other. The rings are conducting with resistance $\lambda$ per unit length but are electrically insulated from each other. This arrangement is kept in such a manner that one ring is horizontal and the other vertical and is supported at diametrically opposite points of one ring at supports $S_1$ and $S_2$. At $t=0$ a magnetic field B is switched on given by $\vec B = B_0\left(\hat i + t^2\hat j\right)$ in the region. Choose the correct alternative(s).
\qfig{p1_phy_s1_q03_1}
\begin{qoptions}
\opt{A}{As time progresses the sum of reaction at $S_1$ and $S_2$ remain constant.}
\opt{B}{Reaction of $S_1$ increases and $S_2$ decreases.}
\opt{C}{Difference of reaction at $S_1$ and $S_2$ at any time t equals $\dfrac{B_0^2\pi R^2 t}{\lambda}$.}
\opt{D}{The time when the ring losses contact from anyone of the contacts is $t=\dfrac{mg\lambda}{B_0^2\pi R^2}$.}
\end{qoptions}

\qn{4} A dielectric sphere of radius R has a spherical cavity of radius $\eta R$ as shown in the figure. Uniformly distributed charge is in the dielectric material of density $\rho$. The center of cavity 'C' is the origin of a cartesian co-ordinate system shown. Points A, B and D are on the cavity surface on the axes. Choose the correct statements of the following (Ignore gravity)
\qfig{p1_phy_s1_q04_1}
\begin{qoptions}
\opt{A}{For $\eta\to 0$, the electric field in the cavity is $\dfrac{\rho R}{3\varepsilon_0}$.}
\opt{B}{A charge particle of mass m and charge q is released from B. For $\eta=\dfrac12$, the time taken by the particle to hit the surface of cavity will be $t=\sqrt{\dfrac{12m\varepsilon_0}{\rho q}}$.}
\opt{C}{The same charge particle mentioned in option (B) is projected from point A towards C with speed $v_1$ for $\eta=\dfrac12$ so that it hits the surface of cavity at a point on x-axis. The speed must be $v_1=\dfrac{R}{2}\sqrt{\dfrac{\rho q}{6m\varepsilon_0}}$.}
\opt{D}{The change in kinetic energy of the charge particle mentioned in option (C) during it's motion before it collides is $\dfrac{\rho q R^2}{6\varepsilon_0}$.}
\end{qoptions}

\qn{5} In a certain arrangement a fine pencil of laser light is made to fall on one of the equal faces AB (AB = BC) of an isosceles prism. The prism is resting on a horizontal surface on its triangular base.
\qfig{p1_phy_s1_q05_1}
The laser beam makes angle $\theta=15^\circ$ with face AC in the horizontal plane. There are three prisms which will be kept in the above arrangement one by one. Prism-1 has a prism angle $A=60^\circ$ and the refractive index of material is $n_1=\sqrt2$. Prism-2 has $A=30^\circ$ and its material has refractive index of $n_2$ and Prism-3 has $A=90^\circ$ and is made up of material of refractive index $n_3$. The emergent rays in the case of each prism is parallel to each other.
\begin{qoptions}
\opt{A}{For prism-1 the emergent ray is deviated by $30^\circ$.}
\opt{B}{For prism-2, $n_2=2\cos\left(\dfrac{\pi}{12}\right)$.}
\opt{C}{For prism-3, $n_3=\sqrt{\dfrac32}$.}
\opt{D}{For all the three prisms the beam is parallel to face AC inside the prism.}
\end{qoptions}

\qn{6} A cycle wheel has a rim of mass equal to the mass of its spokes. The spokes are all identical and symmetrical. This wheel is released on an inclined plane of inclination $\theta$. Friction is sufficient to sustain rolling. Mark the correct statements :- (Take $g=10\ ms^{-2}$)
\qfig{p1_phy_s1_q06_1}
\begin{qoptions}
\opt{A}{Acceleration of the center of wheel is $\dfrac{24}{27}g\sin\theta$.}
\opt{B}{Minimum coefficient of friction necessary for pure rolling is $\dfrac25\tan\theta$.}
\opt{C}{The ratio of kinetic energy of two spokes one at lowest position and one at highest position is $1:7$ at a given moment after rolling starts.}
\opt{D}{For a wheel of mass 5 kg, radius 40 cm and $\theta=30^\circ$, the rate of change of angular momentum about any point on the incline plane is 10 kg m$^2$/s.}
\end{qoptions}

\sect{Section-I (ii): Matching List (Single Correct)}

\qn{7} A uniform ring of mass m = 10 kg and radius R = 5 cm is kept at rest on a smooth horizontal surface in vertical plane as shown in the figure. The ring carries a current i = 7 A. A constant horizontal magnetic field B = 100 mT is switched on at t = 0.
\qfig{p1_phy_s2_q07_1}
List-I describe the statement regarding ring and and List-II gives their respective values.
\begin{center}
\begin{tabular}{|c|p{4.3cm}|c|l|}
\hline
\multicolumn{2}{|c|}{List-I} & \multicolumn{2}{c|}{List-II} \\
\hline
(I) & The magnetic moment of loop in SI unit & (P) & 0.0880 \\
\hline
(II) & Torque acting on loop due to magnetic field in SI unit at t = 0 & (Q) & 0.044 \\
\hline
(III) & Angular acceleration of loop in SI unit at t = 0 & (R) & 0.0055 \\
\hline
(IV) & Kinetic energy of ring when it has rotated through $90^\circ$ & (S) & 0.4400 \\
\hline
 & & (T) & 0.0550 \\
\hline
\end{tabular}
\end{center}
\begin{qoptions}
\opt{A}{I$\to$T; II$\to$R; III$\to$S; IV$\to$R}
\opt{B}{I$\to$P; II$\to$Q; III$\to$R; IV$\to$S}
\opt{C}{I$\to$S; II$\to$T; III$\to$P; IV$\to$Q}
\opt{D}{I$\to$T; II$\to$S; III$\to$P; IV$\to$R}
\end{qoptions}

\qn{8} Three resistors of resistance 6 M$\Omega$, 3 M$\Omega$ and 3 M$\Omega$, one capacitor of capacitance $\dfrac15\,\mu F$, two cells of potential difference 5V and 10V and two switches are connected in the circuit as shown in figure. Both switches closed simultaneously at t = 0 ($e^{-1}=0.37$). List-I describes various electrical quantities in the circuit at particular time and List-II gives the possible values of charges and current.
\qfig{p1_phy_s2_q08_1}
Which one of the following options is correct?
\begin{center}
\begin{tabular}{|c|p{4.3cm}|c|l|}
\hline
\multicolumn{2}{|c|}{List-I} & \multicolumn{2}{c|}{List-II} \\
\hline
(I) & The maximum charge on the capacitor in $\mu C$ & (P) & 0.63 \\
\hline
(II) & Charge on capacitor at t = 1s in $\mu C$ & (Q) & 1.54 \\
\hline
(III) & Current in the branch of capacitor at t = 1s in $\mu A$ & (R) & 2.06 \\
\hline
(IV) & Current passing through battery of 5V at t = 1s in $\mu A$ & (S) & 1.00 \\
\hline
 & & (T) & 0.37 \\
\hline
\end{tabular}
\end{center}
\begin{qoptions}
\opt{A}{I$\to$P; II$\to$T; III$\to$R; IV$\to$S}
\opt{B}{I$\to$T; II$\to$P; III$\to$Q; IV$\to$R}
\opt{C}{I$\to$P; II$\to$T; III$\to$T; IV$\to$Q}
\opt{D}{I$\to$S; II$\to$P; III$\to$T; IV$\to$Q}
\end{qoptions}

\qn{9} A plane mirror is tied to the free end of an ideal spring. The other end of the spring is attached to a wall. The spring with mirror is held vertically to the floor can slide along it smoothly. When the spring is at its natural length, the mirror is found to move at speed v with respect to ground. An object moving towards mirror with speed 2v with respect to ground. Match the following: (Assume no collision before mirror comes to rest for 1st time)
\qfig{p1_phy_s2_q09_1}
\begin{center}
\begin{tabular}{|c|p{4.8cm}|c|l|}
\hline
\multicolumn{2}{|c|}{List-I} & \multicolumn{2}{c|}{List-II} \\
\hline
(I) & Speed of image with respect to ground when spring at natural length. & (P) & v \\
\hline
(II) & Speed of image with respect to mirror when spring at natural length. & (Q) & 0 \\
\hline
(III) & Speed of image with respect to object when spring is at natural length. & (R) & 2v \\
\hline
(IV) & Speed of image with respect to ground frame when spring at maximum compressed state. & (S) & 3v \\
\hline
 & & (T) & 4v \\
\hline
\end{tabular}
\end{center}
\begin{qoptions}
\opt{A}{I$\to$Q; II$\to$R; III$\to$P; IV$\to$S}
\opt{B}{I$\to$T; II$\to$R; III$\to$S; IV$\to$P}
\opt{C}{I$\to$Q; II$\to$P; III$\to$R; IV$\to$R}
\opt{D}{I$\to$T; II$\to$Q; III$\to$R; IV$\to$P}
\end{qoptions}

\qn{10} Consider the position of different nuclei on the binding energy per nucleon graph as shown. List-I has different nuclear reaction involving nuclei A, B, C, D and E. List-II has corresponding Q values in MeV.
\qfig{p1_phy_s2_q10_1}
\begin{center}
\begin{tabular}{|c|l|c|l|}
\hline
\multicolumn{2}{|c|}{List-I} & \multicolumn{2}{c|}{List-II} \\
\hline
(I) & $2B\to C$ & (P) & Zero \\
\hline
(II) & $2B+C\to D$ & (Q) & 102 \\
\hline
(III) & $E\to C+D$ & (R) & Non Zero \\
\hline
(IV) & $3A\to C$ & (S) & 36 \\
\hline
 & & (T) & 90 \\
\hline
\end{tabular}
\end{center}
Which of the following is correct?
\begin{qoptions}
\opt{A}{I$\to$P,S; II$\to$P; III$\to$R,T; IV$\to$Q,R}
\opt{B}{I$\to$R,S; II$\to$P; III$\to$P,T; IV$\to$Q,R}
\opt{C}{I$\to$S,T; II$\to$P,R; III$\to$R,T; IV$\to$Q,R}
\opt{D}{I$\to$R,S; II$\to$P; III$\to$R,T; IV$\to$Q,R}
\end{qoptions}

\sect{Section-II: Numerical Value}

\qn{1} A cylindrical coaxial cable consists of an internal solid conductor of radius 'a' with external thin walled conductor of radius '4a'. The current distribution over the cross-section of the internal conductor is uniform. If inductance per unit length of cable is $\dfrac{n\mu_0}{\pi}$ then find value of n.

\qn{2} A spherical region of radius $R=\sqrt3$ m has uniform charge density $\rho\ C/m^3$. A plane cuts the spherical region at some distance from the center such that the electric flux through the portion of the plane bounded by the sphere is maximum. The maximum flux is found to be $\dfrac{\rho\pi}{N\epsilon_0}$. Find N.

\qn{3} In given arrangement a point object is placed at 36 cm left of equiconvex lens of radius of curvature 90 cm and refractive index $\mu_1$. If final image formed at 20 cm distance from equiconcave lens of refractive index '$\mu_2$' of radius of curvature 60 cm, find product of $\mu_1$ and $\mu_2$. Also given that magnitude of focal length of concave lens is double that of convex lens.
\qfig{p1_phy_s3_q03_1}

\qn{4} A monochromatic light source S is placed 0.5mm over a horizontal plane which has a small mirror in it as shown in figure. A vertical screen is kept 80 cm from the mirror. On the screen an interference pattern is observed. Find the height of the last maxima counting upwards in mm. (Given wave length $\lambda=500$ nm)
\qfig{p1_phy_s3_q04_1}

\qn{5} A small body of mass 'm' is projected up from surface of given planet as shown in figure. If small body has velocity which is just sufficient to get out from the gravity of planet then time taken by small body to reach height of 3R from the surface of planet is $k\sqrt{\dfrac{2R^3}{GM}}$. Find value of k.
\qfig{p1_phy_s3_q05_1}

\qn{6} In a glass bulb filled with gaseous neon two metal plates are used as electrodes separated by a distance d = 4 mm that is much smaller than the linear dimensions of the plates. Ionization energy for neon atom is $E_0=21.5$ eV, average distance travelled by an electron between two consecutive collisions with neon atoms is $\ell=0.4$ mm and charge on an electron is $e=-1.6\times10^{-19}$ C. How much potential difference (in V) between the electrodes should be applied to initiate a discharge process?

\qn{7} A small disk of wood of density $\rho_1=500\ kg/m^3$ floats in a stable orientation in liquid water of density $\rho_2=10^3\ kg/m^3$ and surface tension S = 0.07 N/m. The disk has a diameter of 28mm and a thickness of 10mm. Calculate the depth h in mm that the disk sinks below the free surface considering that the water wets the surface of wood with a contact angle of $0^\circ$. (Take $g=10\ m/s^2$)

\qn{8} In an experiment to measure acceleration of a glider moving down a sloping air track, its velocities $v_1$ and $v_2$ are measured at two points and the time interval between them is also measured and given as $v_1=0.21\pm0.05$, $v_2=0.85\pm0.05$ both in (m/s) and $t=8.0\pm0.1$ s. The acceleration is finally quoted as $a\pm\Delta a$. Find $|a|+|\Delta a|$ in $m/s^2$.

\partsec{2}{CHEMISTRY}
\sect{Section-I (i): One or More Correct}

\qn{1} Consider the following oligosaccharide, formed by the condensation reaction of the monosaccharides, unit I and unit II as shown:
\qfig{p1_chem_s1_q01_1}
Which statements is/are CORRECT:
\begin{qoptions}
\opt{A}{The glycosidic bond is 1, 4-$\alpha$-glycosidic bond.}
\opt{B}{Unit I is an L sugar while unit II is a D sugar.}
\opt{C}{The acyclic form of the monosaccharide UNIT-I is the C5-epimer of D-glucose.}
\opt{D}{The acyclic form of the monosaccharide UNIT-II is the C4-epimer of D-glucose.}
\end{qoptions}

\qn{2} For, $N_2(g)+3H_2(g)\to 2NH_3(g)$; $\Delta H=-99$ kJ/mol at 300K. The molar entropy data of the chemicals at 300K is:
\begin{center}
\begin{tabular}{|c|c|}
\hline
Chemical & Entropy $JK^{-1}mol^{-1}$ \\
\hline
$N_2$ & 190 \\
\hline
$H_2$ & 150 \\
\hline
$NH_3$ & 200 \\
\hline
\end{tabular}
\end{center}
Which relation(s) is/are CORRECT:
\begin{qoptions}
\opt{A}{$\Delta S_{reaction}=-240\ kJ\,K^{-1}mol^{-1}$}
\opt{B}{$\Delta S_{universe}=+90\ J\,K^{-1}mol^{-1}$}
\opt{C}{$\Delta S_{surrounding}=+330\ J\,K^{-1}mol^{-1}$}
\opt{D}{$\Delta S_{reaction}=+240\ J\,K^{-1}mol^{-1}$}
\end{qoptions}

\qn{3} Which of the following is/are correct order of the mentioned property
\begin{qoptions}
\opt{A}{Acidic Strength:\\ \qfigi{p1_chem_s1_q03_1}}
\opt{B}{Dipole moment:\\ \qfigi{p1_chem_s1_q03_2}}
\opt{C}{Heat of hydrogenation:\\ \qfigi{p1_chem_s1_q03_3}}
\opt{D}{Basic strength:\\ \qfigi{p1_chem_s1_q03_4}}
\end{qoptions}

\qn{4} The INCORRECT fact(s) regarding the modern periodic table is/are:
\begin{qoptions}
\opt{A}{The element that has the maximum number of electrons in the last shell belongs to the d-block}
\opt{B}{The atom that has the highest electron affinity is fluorine.}
\opt{C}{Element having electronic configuration \resizebox{\linewidth}{!}{$1s^22s^22p^63s^23p^63d^{10}4s^24p^64d^{10}4f^{14}5s^25p^65d^{10}6s^26p^66d^27s^2$} belongs to the f-block}
\opt{D}{The ionisation energy increases continuously on moving left to right in all the periods.}
\end{qoptions}

\qn{5} At temperature 300 K, a rigid vessel fitted with a valve has He(g) in it at 1 atm and some liquid, A(l) (molar mass = 200 g/mol, vapour pressure = 0.5 atm at 300 K) is placed in a room. Liquid A now vaporises. The valve opens as soon as the pressure inside the vessel becomes more than 1.2 atm and shuts after the pressure gets maintained at or below 1.2 atm.
\qfig{p1_chem_s1_q05_1}
Which of the following is/are INCORRECT for the situation inside the vessel after a very longtime.
\begin{qoptions}
\opt{A}{A stable pressure of 1.2 atm would be developed due to the pressure of He alone.}
\opt{B}{The mole fraction of He in the gaseous phase inside the container would be $\dfrac{5}{12}$.}
\opt{C}{The mole fraction of the vapours of A(l) in the gaseous phase inside the container would be $\dfrac{5}{12}$.}
\opt{D}{The system would never attain a stable equilibrium pressure.}
\end{qoptions}

\qn{6} Which of the following reaction/s is/are correct?
\begin{qoptions}
\opt{A}{$XeO_3+NaOH\to Na[HXeO_4]$}
\opt{B}{$S_8+24XeF_2\to 8SF_6+24Xe$}
\opt{C}{$2XeF_6+SiO_2\to XeOF_4+SiF_4$}
\opt{D}{$XeF_6+6H_2O\to XeO_3+6HF$}
\end{qoptions}

\sect{Section-I (ii): Matching List (Single Correct)}

\qn{7} Match the following:
\begin{center}
\begin{tabular}{|c|p{2.3cm}|c|p{2.3cm}|}
\hline
\multicolumn{2}{|p{2.6cm}|}{\centering List-I (Class of Drugs)} & \multicolumn{2}{p{2.6cm}|}{\centering List-II (Example of the Drugs)} \\
\hline
(P) & Antacid & (1) & Salvarsan \\
\hline
(Q) & Tranquilizer & (2) & Cimetidine \\
\hline
(R) & Analgesic & (3) & Codeine \\
\hline
(S) & Antibiotic & (4) & Meprobamate \\
\hline
 & & (5) & Morphine \\
\hline
\end{tabular}
\end{center}
\begin{qoptions}
\opt{A}{P$\to$1;Q$\to$2;R$\to$3;S$\to$4}
\opt{B}{P$\to$2;Q$\to$3;R$\to$4;S$\to$1}
\opt{C}{P$\to$2;Q$\to$4;R$\to$3;S$\to$1}
\opt{D}{P$\to$1;Q$\to$2;R$\to$4;S$\to$3}
\end{qoptions}

\qn{8} Match the column:
\begin{center}
\begin{tabular}{|c|p{3.6cm}|c|p{2.0cm}|}
\hline
\multicolumn{2}{|p{3.9cm}|}{\centering Column-I\\ The particle under discussion} & \multicolumn{2}{p{2.3cm}|}{\centering Column-II\\ Approximate value of the property associated with the particle} \\
\hline
(P) & The electron in an orbital (with only one angular node) of an isolated H-atom has RPDF (Radial Probability Distribution Function) as shown \newline \qfigi{p1_chem_s2_q08_1} & (1) & Magnitude of the total energy 3.4 eV. \\
\hline
(Q) & An isolated photon of the wavelength = 3647\AA. & (2) & Number of radial nodes in the described orbital = 0. \\
\hline
(R) & The electron of an isolated He$^+$ ion present in the orbital described by $\Psi_{400}$. & (3) & Magnitude of the total energy 1879 eV. \\
\hline
(S) & An isolated electron of the wavelength = 6.64 \AA. & (4) & Number of radial nodes in the described orbital = 3. \\
\hline
 & & (5) & Magnitude of the total energy = 0.2 eV. \\
\hline
\end{tabular}
\end{center}
\begin{qoptions}
\opt{A}{P$\to$1,2;Q$\to$1;R$\to$1,2;S$\to$3}
\opt{B}{P$\to$1,4;Q$\to$3;R$\to$1,3;S$\to$4}
\opt{C}{P$\to$1,5;Q$\to$3;R$\to$1,4;S$\to$4}
\opt{D}{P$\to$1,2;Q$\to$1;R$\to$1,4;S$\to$1}
\end{qoptions}

\qn{9} An acidic indicator HIn ($K_a=10^{-9}$) has the ionised form of red colour and the unionised form of blue colour. The red colour is distinctly visible to a human eye when the ionised form is at least 2 times the unionized form whereas blue colour is distinctly visible to a human eye when the unionised form is at least, 10 times the ionised form. Various titrations are studied using this indicator. When colour change is visible during the titration, end point is observed for different titrations given in list-I. [Given: $\log2=0.3$]
Match the lists:
\begin{center}
\begin{tabular}{|c|p{3.8cm}|c|p{2.0cm}|}
\hline
\multicolumn{2}{|p{4.1cm}|}{\centering List-I (Titration)} & \multicolumn{2}{p{2.3cm}|}{\centering List-II (pH-values)} \\
\hline
(P) & 0.1 M 100 ml HCl (in the flask) vs 0.1 M NaOH (in the burette). & (1) & Equivalence point pH = 9. \\
\hline
(Q) & {\sloppy 0.2 M 100 ml CH$_3$COOH $K_a=10^{-5}$ (in the flask) vs 0.2 M NaOH (in the burette).\par} & (2) & End point pH = 7. \\
\hline
(R) & 0.2 M 100 ml NaOH (in the flask) vs 0.2 M CH$_3$COOH $K_a=10^{-5}$ (in the burette). & (3) & End point pH = 8. \\
\hline
(S) & 0.1 M 100 ml NaOH (in the flask) vs 0.1 M HCl (in the burette). & (4) & End point pH = 9.3. \\
\hline
 & & (5) & End point pH = 10. \\
\hline
\end{tabular}
\end{center}
\begin{qoptions}
\opt{A}{P$\to$4;Q$\to$1,4;R$\to$1,3;S$\to$3}
\opt{B}{P$\to$5;Q$\to$1,5;R$\to$1,3;S$\to$3}
\opt{C}{P$\to$2;Q$\to$2;R$\to$2;S$\to$2}
\opt{D}{P$\to$2;Q$\to$1;R$\to$1;S$\to$2}
\end{qoptions}

\qn{10} Stepwise and overall formation constants ($K_f$) for $[Cu(NH_3)_4]^{2+}$ mentioned as follows:
\resizebox{\linewidth}{!}{$[Cu(H_2O)_4]^{2+}(aq.)+NH_3(aq.)\rightleftharpoons[Cu(NH_3)(H_2O)_3]^{2+}(aq.)+H_2O:K_1$}\par
\resizebox{\linewidth}{!}{$[Cu(NH_3)(H_2O)_3]^{2+}(aq.)+NH_3(aq.)\rightleftharpoons[Cu(NH_3)_2(H_2O)_2]^{2+}(aq.)+H_2O:K_2$}\par
\resizebox{\linewidth}{!}{$[Cu(NH_3)_2(H_2O)_2]^{2+}(aq.)+NH_3(aq.)\rightleftharpoons[Cu(NH_3)_3(H_2O)]^{2+}(aq.)+H_2O:K_3$}\par
\resizebox{\linewidth}{!}{$[Cu(NH_3)_3(H_2O)]^{2+}(aq.)+NH_3(aq.)\rightleftharpoons[Cu(NH_3)_4]^{2+}(aq.)+H_2O:K_4$}\par
\resizebox{\linewidth}{!}{$[Cu(H_2O)_4]^{2+}(aq.)+4NH_3(aq.)\rightleftharpoons[Cu(NH_3)_4]^{2+}(aq.)+4H_2O:K_5$}\par
List-II has the accurate values of the logarithm of the equilibrium constants as described in the above chemical equations. Match the $\log K_f$ in the List-I with their appropriate values given in the List-II.
\begin{center}
\begin{tabular}{|c|c|c|c|}
\hline
\multicolumn{2}{|c|}{List-I} & \multicolumn{2}{c|}{List-II} \\
\hline
(P) & $\log K_1$ & (1) & 2.0 \\
\hline
(Q) & $\log K_2$ & (2) & 2.7 \\
\hline
(R) & $\log K_3$ & (3) & 3.2 \\
\hline
(S) & $\log K_5$ & (4) & 4.0 \\
\hline
 & & (5) & 11.9 \\
\hline
\end{tabular}
\end{center}
\begin{qoptions}
\opt{A}{P$\to$5;Q$\to$4;R$\to$3;S$\to$2}
\opt{B}{P$\to$4;Q$\to$3;R$\to$2;S$\to$5}
\opt{C}{P$\to$2;Q$\to$4;R$\to$3;S$\to$5}
\opt{D}{P$\to$4;Q$\to$3;R$\to$2;S$\to$1}
\end{qoptions}

\sect{Section-II: Numerical Value}

\qn{1} How many of the following statements with regards to the Lanthanides (Ln) is/are correct: \\
(i) The common oxidation state of Ln is +3\\
(ii) $Ce(SO_4)_2$ is coloured\\
(iii) Ln can produce $H_2(g)$ when it reacts with dil. HCl as well as with $H_2O$.\\
(iv) All $Ln(OH)_3$ are basic\\
(v) Mischmetal, an alloy of lanthanoids has approx. 5\% iron and traces of S, C, Ca and Al.

\qn{2} Find the sum total of all the locants position present in the principle chain in the proper IUPAC name of the following compounds:\\
(i) Ethylidene chloride\\
(ii) Mesitylene\\
(iii) Catechol\\
(iv) \qfigi{p1_chem_s3_q02_1}\\
(v) $C_6H_5O(CH_2)_6CH_3$

\qn{3} $C_6H_5CH_3\ (X)\xrightarrow[\ H_2SO_4\ ]{H_2SO_4}A\xrightarrow[\ H_2SO_4\ ]{HNO_3}B\xrightarrow{H_3O^+}C\xrightarrow[\ Cl_2\ ]{FeCl_3}D\xrightarrow[\ h\nu\ ]{2\text{ equivalents } Cl_2(g)}E\xrightarrow[\ H_2O\ ]{Ca(OH)_2}F$\\
$C_6H_5CH_3\ (X)\xrightarrow[AlCl_3]{AcCl}G\xrightarrow[H^+]{Zn/Hg}H\xrightarrow[H^+]{F}I$\\
Consider the above reaction sequence. Find the degree of unsaturation of I.

\qn{4} 100 ml of a dilute $HNO_3$ solution reacts completely with 6.5g powder of pure Zinc. The same sample of dilute $HNO_3$ when used to find the percentage purity of a 10 g impure copper sample needed 80 ml of the acid for the complete reaction. Find the percentage purity of copper in the sample. [Atomic mass: Zn = 65, Cu = 63.5]

\qn{5} Pure $H_2O$ and Sucrose solution are placed separated by an SPM as shown, assuming no change in the temperature throughout the process. Find h (in cm).
\qfig{p1_chem_s3_q05_1}
Use: density of water = density of solution = 1 g ml$^{-1}$, acceleration due to gravity = 10 $ms^{-2}$, and universal gas constant R = 8 J mol$^{-1}K^{-1}$.

\qn{6} The reaction $PCl_5(g)\rightleftharpoons PCl_3(g)+Cl_2(g)$ is at equilibrium. Now suddenly the volume is doubled. The system is now allowed to attain a new equilibrium state. How many of the following are lesser at the new equilibrium as compared to the old equilibrium? Assume the entire process to be isothermal.\\
(a) Moles of $PCl_5$\\
(b) Moles of $PCl_3$\\
(c) Moles of $Cl_2$\\
(d) Concentration of $PCl_5$\\
(e) Concentration of $PCl_3$\\
(f) Concentration of $Cl_2$\\
(g) Rate of the forward reaction\\
(h) Rate of the backward reaction\\
(i) Equilibrium constant

\qn{7} With how many of the following reagents does $H_2O_2$ show oxidising action?\\
(a) $Fe^{2+}$ (Acidic medium)\\
(b) $Mn^{2+}$ (Basic medium)\\
(c) $Mn^{2+}$ (Acidic medium)\\
(d) $I^-$ (Basic medium)\\
(e) $MnO_2$ (Basic medium)

\qn{8} The alkene $C_4H_8$, that has the minimum possible heat of combustion, is subjected to the following sequence of reactions:\\
$C_4H_8\xrightarrow{O_3,\ Zn}A\ \text{(Positive tollen's test)}+B\ \text{(Negative tollen's test)}$\\
$A\xrightarrow{NH_3}C\ \text{(Used to treat urinary infection)}$\\
How many N--N linkage in present in C

\partsec{3}{MATHEMATICS}
\sect{Section-I (i): One or More Correct}

\qn{1} P is a point on $y^2=12x$, such that its focal distance is 6. Tangent drawn at P intersect tangent at vertex at point T(y>0). Point R is on x-axis inside parabola such that SR = 6 (S is focus). M is a point on parabola such that tangents at M \& P meet at point N such that SN = 5, then
\begin{qoptions}
\opt{A}{Area of quadrilateral PRST is 18 units}
\opt{B}{Area of quadrilateral PRST is 27 units.}
\opt{C}{Sum of squares of possible ordinates of M is 28.}
\opt{D}{Sum of squares of possible ordinates of M is 30.}
\end{qoptions}

\qn{2} $f(x)$ is a quadratic function such that $f(0)=1$ and $\displaystyle\int\dfrac{f(x)\cdot dx}{x^2\cdot(x+1)^3}$ is a rational function, where: $\displaystyle\lim_{x\to\infty}\left(\sqrt{f(x)-1}-x\right)=\ell$ (finite) then:
\begin{qoptions}
\opt{A}{$f'(0)=3$}
\opt{B}{$f'(0)=4$}
\opt{C}{leading coefficient of $f(x)$ is 1}
\opt{D}{$\ell=3$}
\end{qoptions}

\qn{3} A triangle ABC has sides AB = 13, BC = 5 and AC = 12 units, slides on XY axes such that A and B remains on positive X and positive Y axis respectively, then: (where [.] is G.I.F.)
\begin{qoptions}
\opt{A}{Locus of C is $5x=12y$}
\opt{B}{Locus of C is $12x=5y$}
\opt{C}{If $p\cdot OA+q\cdot OB=OC$, then $\left[\dfrac1p+\dfrac1q\right]=3$}
\opt{D}{If $p\cdot OA+q\cdot OB=OC$, then $\left[\dfrac1p+\dfrac1q\right]=2$}
\end{qoptions}

\qn{4} If m, n are roots of $2x+\dfrac1x=2$ and $g(x)=\dfrac{8x^6+1}{x^3}$ then
\begin{qoptions}
\opt{A}{$g(m)+g(n)=-8$}
\opt{B}{$g(m)-g(n)=0$}
\opt{C}{$g(m)+g(n)=8$}
\opt{D}{$g(m)-g(n)=2$}
\end{qoptions}

\qn{5} Consider a plane P passing through A(--1, 3, 2), B(7, 5, 10) and C($\lambda$, 3, $\mu$) and a straight line L with positive direction cosines passing through C bisecting AB and makes equal angles with co-ordinate axes. Let $L_1$, be line parallel to L and passing through O (where 'O' is origin). Which one is/are correct?
\begin{qoptions}
\opt{A}{Area of $\triangle ABC$ is $3\sqrt2$}
\opt{B}{Equation of $L_1$ is $\dfrac{x-1}{1}=\dfrac{y-1}{1}=\dfrac{z-1}{1}$}
\opt{C}{$\lambda+\mu$ equals 8.}
\opt{D}{equation to the plane perpendicular to plane P and continuing $L_1$ is $x-2y+z=0$}
\end{qoptions}

\qn{6} The sum of absolute values of x satisfying equation
$$\begin{vmatrix}2^{11} & 2^{11} & -2^{12}-x\\ -2^{11}+x & 2^{12} & -2^{11}\\ 2^{12} & -2^{11}+x & -2^{11}\end{vmatrix}=0$$
is $\alpha$, then
\begin{qoptions}
\opt{A}{$\alpha$ is divisible by 3}
\opt{B}{$\alpha$ is divisible by 8}
\opt{C}{Number of divisors of $\alpha$ = 26}
\opt{D}{$\alpha$ is divisible by 54}
\end{qoptions}

\sect{Section-I (ii): Matching List (Single Correct)}

\qn{7} Match the following:
\begin{center}
\begin{tabular}{|c|p{4.5cm}|c|c|}
\hline
\multicolumn{2}{|c|}{List-I} & \multicolumn{2}{c|}{List-II} \\
\hline
(I) & Q is image of P(2, 4) in bisector line L of 1st quadrant, S is circle with centre Q and radius 4. A circle M is drawn such that it touch line L at (1, 1) and its radical axis with S passes through P then its radius is & (P) & $12\sqrt2$ \\
\hline
(II) & Tangents are drawn from P(K, 0) on $x^2+y^2=50$, meeting y-axis at Q \& R, possible value of K for which area of $\triangle PQR$ is minimum. & (Q) & $\sqrt{50}$ \\
\hline
(III) & If length of major axis of ellipse $2x^2+xy+y^2=5$ is $2k$ and $k^2=\dfrac{10}{a-b}$, then $a^2b$ is: (a is integer) & (R) & $\dfrac{\sqrt{162}}{2}$ \\
\hline
(IV) & Centre of a circle passing through foci of $\dfrac{x^2}{25}-\dfrac{y^2}{24}=1$ and through the points of intersection of tangents at vertices with tangent at point $(5\sqrt2,\ 2\sqrt6)$ is (h, k) then $|2\cdot\sqrt3 k|$ equal & (S) & 10 \\
\hline
 & & (T) & $\sqrt{162}$ \\
\hline
\end{tabular}
\end{center}
\begin{qoptions}
\opt{A}{I$\to$Q;II$\to$R;III$\to$P;IV$\to$T}
\opt{B}{I$\to$R;II$\to$S;III$\to$T;IV$\to$P}
\opt{C}{I$\to$R;II$\to$P;III$\to$T;IV$\to$Q}
\opt{D}{I$\to$Q;II$\to$S;III$\to$S;IV$\to$P}
\end{qoptions}

\qn{8} A box contains 30 tickets, each with a different number from 1 to 30. Four tickets are drawn at random and without replacement. Let p, q, r, s be the number on the ticket drawn.
\begin{center}
\begin{tabular}{|c|p{4.3cm}|c|l|}
\hline
\multicolumn{2}{|c|}{List-I} & \multicolumn{2}{c|}{List-II} \\
\hline
(I) & The probability that $\min(p,q,r,s)<15$ is & (P) & $\dfrac{19}{609}$ \\
\hline
(II) & The probability that the second smallest number is 20 & (Q) & $\dfrac{248}{261}$ \\
\hline
(III) & The probability that the sum of p+q+r+s is even & (R) & $\dfrac{103}{609}$ \\
\hline
(IV) & The probability that product 'pqrs' is even & (S) & $\dfrac{731}{783}$ \\
\hline
 & & (T) & $\dfrac{2951}{5481}$ \\
\hline
\end{tabular}
\end{center}
\begin{qoptions}
\opt{A}{I$\to$P;II$\to$S;III$\to$Q;IV$\to$R}
\opt{B}{I$\to$P;II$\to$S;III$\to$R;IV$\to$Q}
\opt{C}{I$\to$S;II$\to$P;III$\to$T;IV$\to$R}
\opt{D}{I$\to$S;II$\to$P;III$\to$T;IV$\to$Q}
\end{qoptions}

\qn{9} Match the following:
\begin{center}
\begin{tabular}{|c|p{4.2cm}|c|l|}
\hline
\multicolumn{2}{|c|}{List-I} & \multicolumn{2}{c|}{List-II} \\
\hline
(I) & The value of \resizebox{\linewidth}{!}{$\displaystyle\int_0^1\left(\sqrt[2021]{1-x^{2022}}-\sqrt[2022]{1-x^{2021}}\right)dx$} & (P) & 4 \\
\hline
(II) & If $I=\dfrac{\displaystyle\int_0^{51}[x+25]\,dx}{\displaystyle\int_0^{51}\{x+25\}\,dx}$, then $\dfrac{I}{25}$ is equal to (where $[.]\to$ GIF \& $\{.\}\to$ F.P.F.) & (Q) & 0 \\
\hline
(III) & {\sloppy If \resizebox{0.95\linewidth}{!}{$L=\displaystyle\lim_{n\to\infty}\dfrac1n\left(\sin^2\dfrac{\pi}{2n}+\sin^2\dfrac{2\pi}{2n}+\ldots+\sin^2\dfrac{n\pi}{2n}\right)$}, then 2L is equal to\par} & (R) & 2 \\
\hline
(IV) & If $J=\displaystyle\int_0^{\infty}xe^{-2x}\,dx$ then 8J is equal to & (S) & 3 \\
\hline
 & & (T) & 1 \\
\hline
\end{tabular}
\end{center}
\begin{qoptions}
\opt{A}{I$\to$P;II$\to$Q;III$\to$R;IV$\to$S}
\opt{B}{I$\to$Q;II$\to$P;III$\to$R;IV$\to$T}
\opt{C}{I$\to$Q;II$\to$P;III$\to$T;IV$\to$R}
\opt{D}{I$\to$P;II$\to$Q;III$\to$S;IV$\to$R}
\end{qoptions}

\qn{10} Let $t_1=(\sin^{-1}x)^{\sin^{-1}x}$, $t_2=(\sin^{-1}x)^{\cos^{-1}x}$, $t_3=(\cos^{-1}x)^{\sin^{-1}x}$, $t_4=(\cos^{-1}x)^{\cos^{-1}x}$. Match the entries of column-I with the entries of column-II
\begin{center}
\begin{tabular}{|c|c|c|l|}
\hline
\multicolumn{2}{|c|}{List-I} & \multicolumn{2}{c|}{List-II} \\
\hline
(I) & $x\in(0,\cos1)$ & (P) & $t_1>t_2>t_4>t_3$ \\
\hline
(II) & $x\in\left(\cos1,\dfrac{1}{\sqrt2}\right)$ & (Q) & $t_4>t_3>t_1>t_2$ \\
\hline
(III) & $x\in\left(\dfrac{1}{\sqrt2},\sin1\right)$ & (R) & $t_2>t_1>t_4>t_3$ \\
\hline
(IV) & $x\in(\sin1,1)$ & (S) & $t_3>t_4>t_1>t_2$ \\
\hline
 & & (T) & $t_1>t_4>t_3>t_2$ \\
\hline
\end{tabular}
\end{center}
Which of the following is the only correct combination.
\begin{qoptions}
\opt{A}{I$\to$P;II$\to$Q;III$\to$R;IV$\to$S}
\opt{B}{I$\to$Q;II$\to$S;III$\to$R;IV$\to$P}
\opt{C}{I$\to$T;II$\to$S;III$\to$R;IV$\to$P}
\opt{D}{I$\to$Q;II$\to$S;III$\to$T;IV$\to$R}
\end{qoptions}

\sect{Section-II: Numerical Value}

\qn{1} If $y=\left(1+\dfrac1x\right)^x$, then the value of $\dfrac{2\sqrt{y_2(2)+\dfrac18}}{\left|\log\left(\dfrac32\right)-\dfrac13\right|}$ is (where $y_2=\dfrac{d^2y}{dx^2}$)

\qn{2} If $x,y\in(0,2\pi)$, satisfying: $2\sin x\cdot\sin y+3\cos y+6\cos x\cdot\sin y=7$ then the value of $(\tan^2x+\tan^2y)$ is

\qn{3} If orthogonal trajectory of $x^2+y^2-ay=0$ is passing through (4, 2) then the member of orthogonal trajectory is a circle with diameter \_\_\_\_\_\_

\qn{4} Let P(x) be a quadratic polynomial with complex coefficient whose leading coefficient is 1. Suppose the equation $P(P(x))=0$ has 4 distinct solution $x=3,4,\alpha,\beta$. then the sum of all possible values of $(\alpha+\beta)^2=$

\qn{5} Let $z_1$ and $z_2$ be complex number such that $z_1^2-4z_2=16+20i$ and the roots $\alpha$ and $\beta$ of $x^2+z_1x+z_2+m=0$ for some complex number m satisfies $|\alpha-\beta|=2\sqrt7$; if the sum of maximum and minimum value of $|m|$ is equal to $\lambda$, then $\dfrac{\lambda}{5}$ is equal to

\qn{6} If ${}^{30}C_0\,{}^{20}C_{10}+{}^{31}C_1\,{}^{19}C_{10}+{}^{32}C_2\,{}^{18}C_{10}+\ldots+{}^{40}C_{10}\,{}^{10}C_{10}={}^aC_b$, (where a and b are relative co-prime odd numbers) then $\dfrac{a-b}{4}$ is equal to

\qn{7} There are 20 observations $x_1,x_2,\ldots,x_{20}$. If mean of $(x_i+i)\ \forall\ i=1,2,3,\ldots,20$ is 15.5 and $\displaystyle\sum_{i=1}^{20}(x_i-5)^2=100$. If variance of $(x_i-2)\ \forall\ i=1,2,3,..,20$ is $\lambda$, then $\dfrac{\lambda}{4}$ is equal to

\qn{8} If $adjA=\begin{pmatrix}2 & 0 & 0\\ -1 & 2 & 0\\ 3 & 1 & 4\end{pmatrix}$ and $|A|>0$, then value of $\dfrac12\text{tr}(A)$ is equal to (where adj(A), tr(A) and $|A|$ denotes adjoint, trace of matrix A and determinant of A respectively)

\papersec{2}
\partsec{1}{PHYSICS}
\sect{Section-I (i): Single Correct}

\qn{1} Find effective resistance between A and B.
\qfig{p2_phy_s1_q01_1}
\begin{qoptions}
\opt{A}{$\dfrac{7}{15}\Omega$}
\opt{B}{$\dfrac{31}{15}\Omega$}
\opt{C}{$\dfrac{4}{15}\Omega$}
\opt{D}{$\dfrac{6}{15}\Omega$}
\end{qoptions}

\qn{2} A rope of mass m is placed on top of a sphere as shown in figure. The surface of the sphere is frictionless. Find initial acceleration of the rope when released. ($\alpha_2>\alpha_1$)
\qfig{p2_phy_s1_q02_1}
\begin{qoptions}
\opt{A}{$\dfrac{g(\sin\alpha_2-\sin\alpha_1)}{\alpha_1+\alpha_2}$}
\opt{B}{$\dfrac{g(\sin\alpha_2-\sin\alpha_1)}{\alpha_2-\alpha_1}$}
\opt{C}{$\dfrac{g(\cos\alpha_1-\cos\alpha_2)}{\alpha_1+\alpha_2}$}
\opt{D}{$\dfrac{g(\cos\alpha_1-\cos\alpha_2)}{\alpha_2-\alpha_1}$}
\end{qoptions}

\qn{3} A stream of particles each of mass 'm' and moving with velocity 'v' is incident on a small hole of diameter $D_0$. Estimate the diameter of the beam at a distance L from the hole. ($L\gg D_0$)
\begin{qoptions}
\opt{A}{$D_0+\dfrac{2\hbar L}{mD_0v}$}
\opt{B}{$D_0+\dfrac{\hbar L}{mD_0v}$}
\opt{C}{$D_0+\dfrac{\hbar L}{2mD_0v}$}
\opt{D}{None of these}
\end{qoptions}

\qn{4} A droplet of fluid has a donut like shape with inner radius 10 mm and outer radius 20 mm. Radius of curvature of point A is 5 mm in the plain perpendicular to paper. Find radius of curvature of point B in the plain perpendicular to paper.
\qfig{p2_phy_s1_q04_1}
\begin{qoptions}
\opt{A}{10 mm}
\opt{B}{20 mm}
\opt{C}{$\dfrac{20}{3}$ mm}
\opt{D}{$\dfrac{40}{3}$ mm}
\end{qoptions}

\sect{Section-I (ii): One or More Correct}

\qn{5} Monoatomic gas is kept under a piston of mass 3 kg in a vessel in equilibrium. Vessel \& piston are thermally insulated. Piston can move vertically without friction. A block of mass 5 kg is put on top of piston and allowed to attain the equilibrium as in figure (b). The temperature of gas becomes $T_1$. After removing the block piston moves upward and when equilibrium is attained again temperature of gas becomes $T_2$. Neglect atmospheric pressure.
\qfig{p2_phy_s2_q05_1}
\begin{qoptions}
\opt{A}{$T_1=\dfrac53T_0$}
\opt{B}{$T_1=\dfrac73T_0$}
\opt{C}{$T_2=\dfrac54T_0$}
\opt{D}{$T_2=\dfrac74T_0$}
\end{qoptions}

\qn{6} A regular tetrahedron is completely immersed under water of density $\rho$ with it's bottom surface horizontal at a depth of h. Edge length of tetrahedron is 'a'. Neglect atmospheric pressure.
\qfig{p2_phy_s2_q06_1}
\begin{qoptions}
\opt{A}{Net force exerted by the liquid on all the side faces of tetrahedron is $\dfrac{(3\sqrt3h-\sqrt2a)\rho g a^2}{12}$.}
\opt{B}{Force exerted by the liquid on the bottom face of tetrahedron is $\dfrac{\sqrt3\rho g h a^2}{4}$.}
\opt{C}{Net force exerted by the liquid on tetrahedron is $\dfrac{\rho a^3 g}{6\sqrt2}$.}
\opt{D}{Force exerted by liquid on the one of the side face of tetrahedron is $\dfrac{(3\sqrt3h-\sqrt2a)\rho g a^2}{12}$.}
\end{qoptions}

\qn{7} As shown in the figure below, a wheel consists of a large number of thin conducting spokes and is free to rotate about an axle. A brush always makes electrical contact with one spoke at a time at the bottom of the wheel i.e., the circuit never breaks at the brush. A battery with a voltage V flows current through an inductor L into the axle through a spoke, to the brush. A uniform magnetic field B exists perpendicular to the plane of the circuit. The radius of the wheel is R and its moment of inertia about the rotation axis is $I$. At t = 0, the switch is closed. Neglect friction and resistivity of the connecting wires.
\qfig{p2_phy_s2_q07_1}
\begin{qoptions}
\opt{A}{The average kinetic energy of the wheel is $\langle K\rangle=3\dfrac{IV^2}{B^2R^4}$.}
\opt{B}{The current through the inductor is maximum for the first time at $t=\pi\dfrac{\sqrt{IL}}{BR^2}$.}
\opt{C}{The average kinetic energy of the wheel is $\langle K\rangle=\pi\dfrac{IV^2}{B^2R^4}$.}
\opt{D}{The current through the inductor is maximum for the first time at $t=\dfrac{\pi^2}{3}\dfrac{\sqrt{IL}}{BR^2}$.}
\end{qoptions}

\qn{8} A satellite of mass $10^3$ kg is moving in a circular orbit of radius $4\times10^8$ m under gravitational force of a fixed mass at centre O. At a particular point A on the orbit direction of motion of satellite is suddenly changed as shown in figure. As a result motion changes from circular to elliptical. B is nearest point on major axis from O. Distance $OB=8\times10^7$ m. Choose correct options.
\qfig{p2_phy_s2_q08_1}
\begin{qoptions}
\opt{A}{$\dfrac{v_B}{v_A}=6$}
\opt{B}{$\dfrac{v_B}{v_A}=3$}
\opt{C}{eccentricity of new orbit e = 0.8}
\opt{D}{eccentricity of new orbit e = 0.6}
\end{qoptions}

\qn{9} Electric field component of an electromagnetic wave is expressed as $\vec E=\left(3\hat j+b\hat k\right)\times10^{-3}\sin\left[10^7(x+2y+3z-\beta t)\right]$ in SI unit. Taking $c=3\times10^8$ m/s as speed of electromagnetic wave in vacuum. Choose correct options.
\begin{qoptions}
\opt{A}{$\beta=3\sqrt{14}\times10^8$}
\opt{B}{b = 2}
\opt{C}{Average energy density of wave is $6.5\times10^{-6}\varepsilon_0$ in SI unit.}
\opt{D}{Amplitude of magnetic field is $1.2\times10^{-11}$ T.}
\end{qoptions}

\qn{10} A conducting loop of area 0.1 $m^2$ and total resistance 1 $\mu\Omega$ and negligible self inductance is suspended from a non-conducting torsional spring which provides a restoring torque $\tau=-k\phi$, where $\phi$ is the twist angle, in a uniform magnetic field $B=10^{-2}\hat j$. The loop is in the y--z plane at equilibrium and can rotate about z-axis with moment of inertia 0.2 $kgm^2$. The loop is displaced by small angle $\delta$ from equilibrium position and released. Its angular displacement is expressed as $\phi=\delta e^{-\beta t}\left(\cos\gamma t+\dfrac{\beta}{\gamma}\sin\gamma t\right)$. Given k = 3.7 N-m. Choose CORRECT options.
\begin{qoptions}
\opt{A}{Value of $\beta$ is 2.5 $sec^{-1}$.}
\opt{B}{Value of $\beta$ is 3.5 $sec^{-1}$.}
\opt{C}{Value of $\gamma$ is 4.5 $sec^{-1}$.}
\opt{D}{Value of $\gamma$ is 3.5 $sec^{-1}$.}
\end{qoptions}

\sect{Section-III: Single Digit Integer (0-9)}

\qn{1} A charge particle of charge q is released from mid point of two equal fixed charges at a point A. The particle covers a distance AB in time $t_0$. If another charged particle of charge 4q of same mass is released at same point A, it covers same distance AB in time $\dfrac{t_0}{N}$. Find N.

\qn{2} An object of mass 3 kg is moving towards earth of mass M with initial velocity 2 m/s far away from it. Impact parameter is b as shown in figure and b is large enough such that the object does not collide with earth. The final deviation angle is $\theta$ due to gravitational interaction with earth. Consider no other interaction. If $\theta=n\tan^{-1}\left(\dfrac{GM}{mb}\right)$, then find $n\times m$.
\qfig{p2_phy_s3_q02_1}

\qn{3} A object of mass 0.6 kg is released from rest under gravity. Wind is blowing in the horizontal direction with velocity 4 m/s. Due to viscosity of wind the object experience a force $\vec F=-k\vec v$ where $\vec v$ is velocity of object with respect to wind and k = 2 N/m. Find terminal velocity after a long time (in m/s). (Take $g=10\ m/s^2$)
\qfig{p2_phy_s3_q03_1}

\qn{4} A very long cylinder of dust of radius R is spinning about its axis with angular velocity $\omega$. Let r be the distance from axis of the cylinder and $\rho(r)$ it's mass density at distance r. If R = 125 km, $\omega=10^{-7}$ rad/sec, $G=6.7\times10^{-11}\ kg^{-1}m^3\,sec^{-2}$, then find $\dfrac{\rho\left(\frac{3R}{4}\right)}{\rho\left(\frac{R}{2}\right)}$. (Assume mutual gravitational interaction among dust particles.)

\qn{5} A rod of mass M and length L is placed on a rough surface as shown in figure. Friction co-efficient on right side of point P is $2\mu$ and on left side is $\mu$. Linear mass density of rod is $\lambda_0x$ where $\lambda_0$ is constant and x is distance from left end. The rod is moved along the surface completely to the other side where friction co-efficient is $2\mu$. Work done during this process is $\dfrac{nMgL\mu}{3}$. Find n.
\qfig{p2_phy_s3_q05_1}

\qn{6} A parallel beam of light falls perpendicular at one side of the glass semi-toroid which has rectangular cross-section. Refractive index of glass is $\mu=1.5$. The inner radius of glass toroid is r = 10 cm. Find maximum glass thickness d (in cm) when all light entering through cross-section A will emerge from other flat cross-section B.
\qfig{p2_phy_s3_q06_1}

\qn{7} A segment of ring AB can oscillate as in figure in the plane of the ring with the help of string OA and OB. Radius of segment of ring is $R=\dfrac{3}{10\pi}$ m. Length of OA and OB is also equal to R and angle $\phi=\pi/3$. Strings OA and OB are massless. Mass of segment AB is m. Axis of rotation of the segment is normal to paper. Time period of small oscillation is $T=\dfrac{\pi}{N}$ sec. Find N.
\qfig{p2_phy_s3_q07_1}

\qn{8} A horizontal closed organ pipe has right end open to atmosphere and has a very heavy piston, which serves as a closed end. The piston can be slided on either side to adjust the length of air column as shown in the figure. Ignore end correction. It is vibrating in $1^{st}$ overtone with tuning fork of frequency 'f' such that the length of the air column is L = 96 cm. Now tuning fork is replaced by another tuning fork of frequency $\dfrac{32}{15}f$. Find the minimum distance (in cm) by which the piston should be shifted so that standing waves are again formed on the pipe.
\qfig{p2_phy_s3_q08_1}

\partsec{2}{CHEMISTRY}
\sect{Section-I (i): Single Correct}

\qn{1} A unit cell of 3D-HCP is shown here:
\qfig{p2_chem_s1_q01_1}
Let, $x=$ Distance between the center of the Tetrahedral void formed by ABGO and the plane ABCDEF and $y=$ Distance between the center of the Octahedral void formed by ABGOPR and the plane ABCDEF. Then, which of the following is/are true?
\begin{qoptions}
\opt{A}{$x=y=R\sqrt{\dfrac23}$}
\opt{B}{$x<y$}
\opt{C}{$x+y=2R\sqrt{\dfrac23}$}
\opt{D}{$y-x=\dfrac{R}{2}\sqrt{\dfrac23}$}
\end{qoptions}

\qn{2} 'S' is mixture of two salts. Consider the following facts about this mixture:\\
a. Aqueous solution of S $\xrightarrow{H_2SO_4}$ White precipitate (W) + Coloured Solution (C)\\
b. W + $CH_3COONH_4$ (conc) $\to$ Soluble complex\\
Appropriate option(s) for the mixture of the salt(s) is/are:
\begin{qoptions}
\opt{A}{$AgNO_3$ and $Zn(NO_3)_2$}
\opt{B}{$Pb(NO_3)_2$ and $Cu(NO_3)_2$}
\opt{C}{$Zn(NO_3)_2$ and $Cu(NO_3)_2$}
\opt{D}{$Pb(NO_3)_2$ and $Zn(NO_3)_2$}
\end{qoptions}

\qn{3} Which of the following statement(s) is/are CORRECT?\\
(I) Sodium is found as a native ore in the nature.\\
(II) Sodium react with excess oxygen to form $Na_2O_2$.\\
(III) $Na_2O_2$ is paramagnetic in nature.\\
(IV) Sodium reacts with $N_2$ (from air) to form $Na_3N$.
\begin{qoptions}
\opt{A}{II Only}
\opt{B}{II, III, IV}
\opt{C}{II and IV}
\opt{D}{III and IV}
\end{qoptions}

\qn{4} Based on the given reaction sequence, choose the CORRECT statement(s):
\begin{center}\resizebox{\linewidth}{!}{$A\ (C_8H_8O_2)\xrightarrow[\Delta]{NH_3}B\xrightarrow[KOH]{Br_2}C\xrightarrow[\text{ii.}Cu_2(CN)_2,HCN\text{ iii.}H_3O^+]{\text{i.}NaNO_2+HCl}D$}\end{center}
Given, A is the most acidic among all of its isomeric forms.
\begin{qoptions}
\opt{A}{C is the most basic among all of its isomers.}
\opt{B}{A and D both give effervescence with $NaHCO_3$ (aq.).}
\opt{C}{C reacts with the Hinsberg's reagent and forms a compound which does not dissolve in alkali solution.}
\opt{D}{B is more basic than C.}
\end{qoptions}

\sect{Section-I (ii): One or More Correct}

\qn{5} Which is/are CORRECT statement(s):
\begin{qoptions}
\opt{A}{In Kjeldahl's method of estimation of Nitrogen, $CuSO_4$ is added in the Kjeldahl's flask to increase the boiling point of the reaction mixture.}
\opt{B}{Duma's method of estimation of Nitrogen can not be applied to pyridine.}
\opt{C}{In the estimation of oxygen by successively heating the compound strongly and then passing the gas over red hot coke and later through warm $I_2O_5$, each mole of $O_2$ gas liberated from the organic compound will produce two moles of carbon dioxide.}
\opt{D}{Presently the elements carbon, hydrogen and nitrogen are determined by CHN elemental analyzer.}
\end{qoptions}

\qn{6} Under different reaction conditions $Ni^{2+}$ reacts with (P) (excess) $Cl^-$ and (Q) (excess) $CN^-$. Choose CORRECT statement(s) regarding above statement:
\begin{qoptions}
\opt{A}{P \& Q can result in tetrahedral and square planar geometries around nickel respectively.}
\opt{B}{P and Q both form homoleptic complexes}
\opt{C}{When Ligand Q reacts with $Ni^{2+}$ then low spin complex is formed}
\opt{D}{When ligand P reacts with $Ni^{2+}$, diamagnetic complex is formed}
\end{qoptions}

\qn{7} When excess $KMnO_4$ is added to concentrated $H_2SO_4$ an oily green color covalent compound X is formed. Which of the following is/are true?
\begin{qoptions}
\opt{A}{X is formed by a dehydration reaction}
\opt{B}{In X, Mn is octahedrally surrounded by 'O' atoms}
\opt{C}{X is an amphoteric oxide}
\opt{D}{X contains a bridge bond.}
\end{qoptions}

\qfullbegin
\qn{8} Consider the given reactions sequence:
$$(P)\xrightarrow[\text{2.}H_3O^+]{\text{1.}KMnO_4/OH^-/H_2O}(Q)\xrightarrow{\Delta}(R)\xrightarrow[\text{3.}H_3O^+]{\text{1. toluene},AlCl_3\text{; 2.}N_2H_4/OH^-,\Delta}(S)$$
$$(S)\xrightarrow[\text{2.}AlCl_3,\Delta]{\text{1. }SOCl_2}(T)\rightleftharpoons(T')\xrightarrow{CH_3COCl,\ Py}(U)$$
\qfig{p2_chem_s2_q08_1}
\qfullsplit
Choose the CORRECT option(s):
\begin{qoptions}
\opt{A}{U is \qfigi{p2_chem_s2_q08_2}}
\opt{B}{U is \qfigi{p2_chem_s2_q08_3}}
\opt{C}{Q is a monomer of Glyptal}
\opt{D}{R has ROU 7}
\end{qoptions}
\qfullend

\qn{9} P (The heavier monomer of Buna--S) $\xrightarrow[\text{(ii)}H_3O^+]{\text{(i)}KMnO_4,OH^-}Q\xrightarrow[CHCl_3]{SOCl_2}R\xrightarrow{(CH_3)_2Cd}S\xrightarrow{MCPBA}T$
\qfig{p2_chem_s2_q09_1}
Based on the above sequence of chemical reaction which is/are CORRECT?
\begin{qoptions}
\opt{A}{P has DU = 5}
\opt{B}{S is \qfigi{p2_chem_s2_q09_2}}
\opt{C}{T is \qfigi{p2_chem_s2_q09_3}}
\opt{D}{S can undergo aldol condensation}
\end{qoptions}

\qn{10} The density of $SO_2$ gas at 8 atm and 400K is found to be 10 g/L. Which of the following statement(s) is/are true? [Given: R = 0.08 L-atm $mol^{-1}K^{-1}$, Atomic mass: S = 32]
\begin{qoptions}
\opt{A}{The gas shows negative deviation from the ideal behaviour}
\opt{B}{The repulsive forces are dominant in the gas under the given conditions.}
\opt{C}{The compressibility factor at 9 atm, 400K will definitely be more than the compressibility factor at 8 atm, 400K}
\opt{D}{The compressibility factor, Z at 8 atm, 400K is 0.625}
\end{qoptions}

\sect{Section-III: Single Digit Integer (0-9)}

\qn{1} At 300K, the stopping potential for the photoelectrons being emitted from a metal surface of work function 5.4 eV is equal to the standard electrode potential of the cell;\\
Ag(s) $\mid$ AgBr(s) $\mid$ KBr (aq) $\parallel$ KCl (aq) $\mid$AgCl (s) $\mid$Ag(s)\\
When the wavelength of the incident photons on the same metal plate is halved, the wavelength of the most energetic photoelectrons emitted is $x$\AA. Find x (rounded off to the nearest integer)\\
Given $K_{SP}\,AgCl=10^{-10}$, $K_{SP}\,AgBr=10^{-15}$ and $\dfrac{2.303R}{F}=\dfrac{0.06}{300}$

\qn{2} How many of the following have permanent dipole moment and at least one d-orbital is participating in bonding?
\begin{center}
\begin{tabular}{|c|}
\hline
(i)\quad $XeF_4$ \\
\hline
(ii)\quad $XeO_2F_2$ \\
\hline
(iii)\quad $H_2O$ \\
\hline
(iv)\quad $PF_2Cl_3$ \\
\hline
(v)\quad $H_5IO_6$ \\
\hline
(vi)\quad $NO_2$ \\
\hline
(vii)\quad $SF_6$ \\
\hline
(viii)\quad $XeOF_4$ \\
\hline
(ix)\quad $SOF_4$ \\
\hline
\end{tabular}
\end{center}

\qn{3} A sequence of reactions of phosphorous is given below:
\qfig{p2_chem_s3_q03_1}
Sum of the oxidation numbers of the central atoms in the compounds Q,R,S, and T is 3x. The value of x is?

\qn{4} A cube of 1 cm side when cut into $10^{12}$ equal cubes, The total surface area (in $m^2$) of all the cubes thus obtained is?

\qn{5} Find the change in the entropy (in $JK^{-1}$) of the universe when 1 mole $CO_2$ gas (assume ideal behavior) at 6 atm undergoes isothermal expansion against vacuum upto 2 atm (Use R = 8 J/mol$^{-1}K^{-1}$, $\ln3=1$)

\qn{6} 100 ml of a river water sample was treated with 50 ml $K_2Cr_2O_7$ solution, the unreacted $K_2Cr_2O_7$ required 20 ml of 0.1 N Mohr's salt solution, during the back-titration process. In another experiment reaction between the same quality and quantity of the $K_2Cr_2O_7$ solution and 0.1 N Mohr's salt solution, 21 ml of Mohr's salt solution is consumed. The COD of the river water sample is?

\qn{7} $CH_2=CH-(CH_2)_3-CH=CH_2\xrightarrow[\text{(i)}H_2O\text{ (ii)}H^+/\Delta]{Br_2}C_7H_{12}Br_2O\ (P)$\\
Number of stereoisomers of P are?

\qn{8} For the reaction:\\
$KMnO_4+H_2C_2O_4+H_2SO_4\to K_2SO_4+MnSO_4+CO_2+H_2O$\\
The n-factor for $H_2SO_4$ is $\dfrac23x$. The value of x is?

\partsec{3}{MATHEMATICS}
\sect{Section-I (i): Single Correct}

\qn{1} Let $F_n(x)=\displaystyle\int\sin(n+2)x\,\sin^nx\,dx$ and $F_n(\pi)=0$ where $n\in N$. If the value of $F_{16}\left(\dfrac{3\pi}{2}\right)$ is a, then determine $\left(\dfrac1a+2\right)$
\begin{qoptions}
\opt{A}{19}
\opt{B}{25}
\opt{C}{30}
\opt{D}{28}
\end{qoptions}

\qn{2} If the global maximum value of $f(x)=-x^2+px-1-(p+1)^2$ for $x\in[0,1]$ be $-5$, then the sum of different values of $[P]$ is (where [.] is greatest integral function)
\begin{qoptions}
\opt{A}{3}
\opt{B}{2}
\opt{C}{$-2$}
\opt{D}{$-3$}
\end{qoptions}

\qn{3} Let P(x) = 0 be a fifth degree polynomial equation with integer coefficients that has atleast one integral root. If P(2) = 13 and P(10) = 5, then the integral value of 'x' that must satisfy P(x) = 0.
\begin{qoptions}
\opt{A}{17}
\opt{B}{15}
\opt{C}{21}
\opt{D}{26}
\end{qoptions}

\qn{4} The equation of plane which is equally inclined to the lines $\dfrac{x-1}{2}=\dfrac{y}{-2}=\dfrac{z+2}{-1}$ and $\dfrac{x+3}{8}=\dfrac{y-4}{1}=\dfrac{z}{-4}$ and passing through the origin, is
\begin{qoptions}
\opt{A}{$14x-5y-7z=0$}
\opt{B}{$2x+7y+z=0$}
\opt{C}{$3x-4y-z=0$}
\opt{D}{$x+2y-5z=0$}
\end{qoptions}

\sect{Section-I (ii): One or More Correct}

\qn{5} Let $z_k=\cos\dfrac{2k\pi}{11}+i\sin\dfrac{2k\pi}{11}$ (where k = 0, 1, 2, ...., 10). If $z=\dfrac{(z_1+z_2)(z_2+z_3)(z_3+z_4)\ldots(z_{10}+z_{11})(z_{11}+z_1)}{z_1z_2z_3\ldots z_{11}}$, then z may be equal to
\begin{qoptions}
\opt{A}{$\lambda(i^i)$ (For some real value of $\lambda$)}
\opt{B}{purely imaginary}
\opt{C}{purely real}
\opt{D}{None of these}
\end{qoptions}

\qn{6} $2p^3-(2x+4\sin x-\cos x)p^2-(x\cdot\cos x-4x\cdot\sin x+\sin2x)p+x\cdot\sin2x=0$ is a differential equation with a possible solution which is: $\left(p=\dfrac{dy}{dx}\right)$
\begin{qoptions}
\opt{A}{$y=\dfrac{x^2}{2}+1$}
\opt{B}{$y=2\cos x+3$}
\opt{C}{$y=\dfrac{\sin x}{2}+3$}
\opt{D}{$2y+\sin x=5$}
\end{qoptions}

\qn{7} Let $p=\displaystyle\sum_{r=1}^{2021}(-1)^{r-1}\dfrac{{}^{2021}C_r}{2^{r+1}-1}$ and $q=\displaystyle\sum_{r=1}^{\infty}\dfrac{(2^r-1)^{2021}}{2^{2022r}}$ then p + q is
\begin{qoptions}
\opt{A}{a prime number}
\opt{B}{1}
\opt{C}{an odd number}
\opt{D}{6}
\end{qoptions}

\qn{8} Let $f:R\to(0,1)$ be a continuous function then which of the following pair of vectors are linearly dependent for given interval
\begin{qoptions}
\opt{A}{$\vec a=f(x)\hat i+3\hat j$, $\vec b=(x^2+x)\hat i+x\hat j$, $\forall x\in(-1,\ 0)$}
\opt{B}{$\vec a=\left(\displaystyle\int_0^{1-x}f(t)\,dt\right)\hat i+3\hat j$, $\vec b=x\hat i+2\hat j$, $\forall x\in(0,\ 1)$}
\opt{C}{$\vec a=\left(\displaystyle\int_0^{1-x}f(t)\,dt\right)\hat i+2\hat j$, $\vec b=x\hat i+3\hat j$, $\forall x\in(0,\ 1)$}
\opt{D}{$\vec a=f(x)\hat i+4\hat j$, $\vec b=x^2\hat i+3\hat j$, $\forall x\in(0,\ 1)$}
\end{qoptions}

\qn{9} Which of the following is/are correct?
\begin{qoptions}
\opt{A}{If $\alpha,\beta,\gamma$ are roots of $x^3+px^2+q=0$, then $\begin{vmatrix}2\alpha & 3\beta & 4\gamma\\ 2\beta & 3\gamma & 4\alpha\\ 2\gamma & 3\alpha & 4\beta\end{vmatrix}$ is equal to $24p^3$}
\opt{B}{If $3A=\begin{pmatrix}1 & 2 & 2\\ 2 & 1 & -2\\ x & 2 & y\end{pmatrix}$ and $AA'=I$ then $x+y$ is equal to $-3$ (where I is identity matrix order 3)}
\opt{C}{A is orthogonal matrix of odd order if $|A|>0$ then $|A-I|=0$.}
\opt{D}{If A is sq. matrix of order 3, then $\left|(A-A')^{2023}\right|$ is 0}
\end{qoptions}

\qn{10} If $y=mx+c$ be a tangent hyperbola $\dfrac{x^2}{\lambda^2}-\dfrac{y^2}{(\lambda^3+\lambda^2+\lambda)^2}=1$, then value of $16m^2$ can be
\begin{qoptions}
\opt{A}{9}
\opt{B}{12}
\opt{C}{6}
\opt{D}{15}
\end{qoptions}

\sect{Section-III: Single Digit Integer (0-9)}

\qn{1} Let $f(x)=1-|x-2|$ for $1\le x\le3$ and $f(3x)=3f(x)$ for all positive real x. The number of points in [1, 100] where f(x) is non differentiable, is

\qn{2} For $x,y,z,t\in R$, if $\sin^{-1}x+\cos^{-1}y+\sec^{-1}z\ge t^2-\sqrt{2\pi}\,t+3\pi$, then $\sec\left(\tan^{-1}x+\tan^{-1}y+\tan^{-1}z+\tan^{-1}\left(\sqrt{\dfrac2\pi}t\right)\right)$

\qn{3} Let A be the minimum value of the function $f:R\to R$, $f(x)=\dfrac{(x^2-x+1)^3}{x^6-x^3+1}$, then value of $A\left(\dfrac{2+\sqrt3}{\sqrt3}\right)$ is -

\qn{4} In a 'n' sided polygon, if the number of hexagon that can be constructed by joining these vertices such that none of the side of polygon is also the side of hexagon is 49, then $\dfrac{n}{2}$ is equal to

\qn{5} If the inequality, $\sin x+\sin2x+\sin3x+\ldots+\sin nx\le\dfrac{\sqrt3}{2}$ is valid for exhaustive set of $x\in(\ell_1,\ell_2)$ (where $x\in(0,2\pi)$) for some values of n, $n\in N$, then $\dfrac{3|\ell_1-\ell_2|}{2\pi}$ equals

\qn{6} If a, b, c, d, e are positive real numbers such that $a+b+c+d+e=15$ and $ab^2c^3d^4e^5=(120)^3(50)$ then value of $\dfrac{a^2+b^2+c^2+d^2+e^2}{11}$ is equal to

\qn{7} If the circle $C_1$ touches x-axis and the line $y=x\tan\theta$, $\theta\in\left(0,\dfrac{\pi}{2}\right)$ in first quadrant and circle $C_2$ touches the line $y=x\tan\theta$, y-axis and circle $C_1$ in such a way that ratio of radius of $C_1$ to radius of $C_2$ is 2 : 1. If the value of $\tan\dfrac{\theta}{2}=\dfrac{\sqrt a-b}{c}$ where a,b,c are relatively prime natural numbers, then the value of $\dfrac{a+b+c}{11}$ is

\qn{8} If the probability that $(a^2+b^2+c^2)$ is divisible by 7 is $\lambda$ then $21\lambda$ equals (Given $a,b,c\in N$)

\end{multicols}
\end{document}
